\documentclass[10pt,a4paper]{amsart}
\usepackage[margin=19mm]{geometry}
\usepackage{amsmath,amssymb,mathtools}
\usepackage[scr=rsfs,cal=euler]{mathalfa}
\usepackage{xfrac}
\usepackage[unicode,hidelinks,bookmarks=false]{hyperref}
\newcommand{\ie}{i.e.\ }
\newcommand{\cf}{cf.\ }
\newcommand{\resp}{respectively\ }
\newcommand{\Subsection}{\subsection}
\providecommand{\nobreakdash}{\nobreak\hbox{-}}
\setlength{\emergencystretch}{2em}
\multlinegap0pt
\usepackage{amssymb}
\usepackage{enumerate}
\newtheorem{thm}{Theorem}
\title{On the pre-Schwarzian norm estimate of special close-to-convex harmonic mappings}
\author{Sushil Pandit}
\date{Source: arXiv:2606.19022v1 (CC BY 4.0)}
\begin{document}
\maketitle

\noindent\emph{Source and licence.} On the pre-Schwarzian norm estimate of special close-to-convex harmonic mappings. Version arXiv:2606.19022v1 (CC BY 4.0), distributed by arXiv under the Creative Commons Attribution 4.0 International licence, http://creativecommons.org/licenses/by/4.0/, as recorded in the arXiv licence metadata for this exact version and shown on the abstract page https://arxiv.org/abs/2606.19022v1. Full licence notice: \texttt{licenses/arxiv-2606.19022-CC-BY-4.0.txt}.\par\smallskip

\noindent\textit{Editorial scope.} Counted unit: the article's thm p807 (source0.tex lines 259--261). The article's complete original proof of it is delivered with it (source0.tex lines 262--301, one \texttt{proof} environment of 40 lines). No mathematics is written, repaired or re-derived: the statement and the proof are the article's own lines, in the article's own notation, and no retained line is substituted or re-flowed.  Retained uncounted as the source-local context the proof needs: lines 190--192.  Nothing was omitted from any retained span, so the package carries no omission record.  No retained line contains a citation, so the wrapper carries no bibliography and no bibliography entry.  No retained line contains a figure environment, an image-inclusion command or drawing code, so no image is delivered and the package declares no asset dependency.  Editorial formatting only, in addition: an AMS \texttt{amsart} preamble that restates the article's own declarations and macros used by the retained lines, namely the environments \texttt{thm} on separate counters, together with the standard packages those lines need.  The retained environments print in this document as Theorem 1 in order of appearance, whereas the article prints the same items with the numbering of its own full document.  The complete native source of the article as distributed in the arXiv e-print for this exact version is bundled unchanged as \texttt{sources/203111/source0.tex}.
    \begin{align}\label{P-151}
		h(z)=z+\sum\limits_{n=2}^\infty a_nz^n,~g(z)=\sum\limits_{n=1}^\infty b_nz^n.
	\end{align}

	\begin{thm}\label{p807}
		Let $f=h+\overline{g}\in\mathcal{H}_c, -1/2\le c\le 0$ be a harmonic mapping of the form \eqref{P-151}. Then pre-Schwarzian norm $\|P_f\|\le 4(1-c)+1$. The estimate is sharp. 
	\end{thm}

	\begin{proof}
		As $f=h+\overline{g}\in\mathcal{H}_c$, from the definition, it follows that
		$$1+\frac{zh''(z)}{h'(z)}\prec\frac{1+(1-2c)z}{1-z},$$ which implies that,
		$$\frac{zh''(z)}{h'(z)}\prec\frac{2(1-c)z}{1-z}.$$ Thus, we get
		\begin{align}\label{p805}
			\left|\frac{zh''(z)}{h'(z)}\right|\le \frac{2|z|(1-c)}{1-|z|}.
		\end{align}
		This estimate is sharp for the analytic function $h(z)=\frac{(1-z)^{2c-1}-1}{1-2c}.$\\
		
		The pre-Schwarzian derivative $P_f$ of a harmonic mapping $f=h+\overline{g}\in\mathcal{H}_c$ is 
		$$
			P_f=\frac{h''}{h'}-\frac{\overline{\omega}\omega'}{1-|\omega|^2}.
		$$
		Thus from Schwarz Pick Lemma and the relation \eqref{p805}, the pre-Schwarzian norm $\|P_f\|$ is 
		\begin{align}\label{HM-100}
			\|P_f\|= & \sup_{z \in \mathbb{D}}\left|\frac{h''}{h'}-\frac{\overline{\omega}\omega'}{1-|\omega|^2}\right|(1-|z|^2)\\\nonumber
			\le & \sup_{z \in \mathbb{D}}\left(\left|\frac{h''}{h'}\right|+\left|\frac{\overline{\omega}\omega'}{1-|\omega|^2}\right|\right)(1-|z|^2)\\\nonumber
			\le & \sup_{z \in \mathbb{D}}\left(\left|\frac{2(1-c)}{1-|z|}\right|+\frac{1}{1-|z|^2}\right)(1-|z|^2)\\\nonumber
			=& \sup_{z \in \mathbb{D}}\left(2(1+|z|)(1-c)+1\right)\\\nonumber
			=& 4(1-c)+1.
		\end{align}
		
		To show that the estimate is sharp, we consider a harmonic mapping $f_d=h+\overline{g}$ with $h(z)=\frac{(1-z)^{2c-1}-1}{1-2c}$ and dilatation $\omega_d(z)=\frac{z-d}{1-dz},~d\in(0,1).$ It is clear that $f_d\in \mathcal{H}_c$ and so $\|P_{f_d}\|\le 4(1-c)+1.$ A few calculations show that 
		\begin{align}
			P_h(z)=\frac{2(1-c)}{1-z}\quad\text{and}\quad \frac{\overline{\omega_d(z)}\omega_d'(z)}{1-|\omega_d(z)|^2}=\frac{\overline{z}-d}{(1-dz)(1-|z|^2)}
		\end{align}
		from which it follows that 
		\begin{align}
			\|P_{f_d}\|=&\sup_{z \in \mathbb{D}}\left|\frac{h''(z)}{h'(z)}-\frac{\overline{\omega_d(z)}\omega_d'(z)}{1-|\omega_d(z)|^2}\right|(1-|z|^2)\\\nonumber
			=&\sup_{z \in \mathbb{D}}\left|\frac{2(1-c)}{1-z}-\frac{\overline{z}-d}{(1-dz)(1-|z|^2)}\right|(1-|z|^2)\\\nonumber
			\ge & L_d
		\end{align}
		where 
		\begin{align}
			L_d=&\sup_{z \in [0,1)}\left|\frac{2(1-c)}{1-z}-\frac{\overline{z}-d}{(1-dz)(1-|z|^2)}\right|(1-|z|^2)\\\nonumber\\\nonumber
			=&\sup_{r \in [0,1)}\left|2(1-c)(1+r)-\frac{r-d}{1-dr}\right|\\\nonumber
			=&\sup_{r \in [0,1)}M(r)\\\nonumber
		\end{align}
		with $M(r)=2(1-c)(1+r)-\frac{r-d}{1-dr}.$ We note that $M(0)=2(1-c)+d,~\phi(1^-)=4(1-c)-1.$ Next $M'(r)=0$ implies that $$r=\frac{1}{d}\pm\frac{1}{d}\sqrt{\frac{1-d^2}{2(1-c)}}.$$ We consider $\frac{1}{d}-\frac{1}{d}\sqrt{\frac{1-d^2}{2(1-c)}}=r_0$ as it lies in $(0,1).$ Since $r_0$ is stationary point and $$M(r_0)=2(1-c)+\frac{3-2c}{d}-\frac{4(1-c)}{d}\sqrt{\frac{1-d^2}{2(1-c)}}\rightarrow 4(1-c)+1$$ as $d$ tends to $1,$ it follows that $L_d=4(1-c)+1.$ Thus we get the relation $$4(1-c)+1=L_d\le \|P_{f_d}\|\le 4(1-c)+1$$ which implies that $\|P_{f_d}\|=4(1-c)+1.$
	\end{proof}

\end{document}
